% Pursuing Stacks (cotes 134-2 to 134-8): archive cote 134-6, site page 76 (PDF page 77), typescript page 394.
% typescript. Draft compared: ps5.tex:5278-5361.
% AI-transcribed by Claude Opus 5.5 (claude-opus-5-5, adaptive thinking, effort high), Message Batches API, batch ps3 (msgbatch_01XdFmuVVBAcgc7EFxdtCqsj), 2026-11-12; not yet checked by eye.
\scanpage{134-6/76}
\mspage{394}

1\unsure{2}.8. \hand{and 13.8.}

105.    Let's go on with the overall review of abelianization.

D)  \uline{Duality for topoi} \hand{\uline{of the type $A^\wedge$}, \uline{and} \uline{tentative} \uline{generalizations}.}

The main fact it seems which will give rise to duality statements
for topoi of the type $A^\wedge$ is the following, rather familiar one:

\uline{Proposition} 3.   Let A be a small category, \hand{$\mathcal{M}$} a \uline{U}-category \xout{such that}
stable under small direct limits,
\[ (32) \qquad \hand{$\varepsilon$}_A : A \hand{$\hookrightarrow$} A^\wedge \]
the canonical inclusion functor.  Then the following functor is an
equivalence of categories:
\[ (3\over{2}{3}) \qquad F \hand{$\longmapsto$} F \circ \hand{$\varepsilon$}_A : \quad \underline{\mathrm{Hom}}_{\hand{$!$}}(A^\wedge,\hand{$\mathcal{M}$}) \hand{$\longrightarrow$} \underline{\mathrm{Hom}}(A,\hand{$\mathcal{M}$}) \quad . \]

As I am at a loss to give a reference for this standard fact of cate-%
gory theory,  I'll give in guise of a proof the indication that a quasi-
inverse functor for (33) is given by the familiar construction
\[ (34) \qquad i \hand{$\longmapsto$} i_{\hand{$!$}} : \quad \underline{\mathrm{Hom}}(A,\hand{$\mathcal{M}$}) \hand{$\longrightarrow$} \underline{\mathrm{Hom}}_{\hand{$!$}}(A^\wedge,\hand{$\mathcal{M}$}) \quad , \]
where for any functor
\[ i: A \hand{$\longrightarrow$} \hand{$\mathcal{M}$} \quad , \]
the functor  \xout{i\unsure{x}}\hand{$_!$}
\[ i_{\hand{$!$}}: A^\wedge \hand{$\longrightarrow$} \hand{$\mathcal{M}$} \]
is defined by the formula
\[ (35) \qquad i_{\hand{$!$}}(F) = \hand{$\varinjlim$}_{a \text{ in } A_{/F}} \; i(a) \qquad . \]
It is readily checked (and we have already used a number of times)  that
this functor admits a right adjoint
\[ (36) \qquad i^{*} : \hand{$\mathcal{M}$} \hand{$\longrightarrow$} A^\wedge \quad , \quad i^{*}(x) = ( a \hand{$\longmapsto$} \mathrm{Hom}(i(a),x)) \quad , \]
and hence $i_{\hand{$!$}}$ commutes to small direct limits i.e. is in $\underline{\mathrm{Hom}}_{\hand{$!$}}(A^{\wedge},\hand{$\mathcal{M}$})$ ,
hence (34).

\uline{Remark}. If \hand{$\mathcal{M}$} is a pseudocotopos,  then by \inserted{(}\inserted[typed]{prop.3 above and}\inserted{by\}} coro\unsure{l}lary 2 of prop.1 (p.385)
the functor
\[ (37) \qquad i \hand{$\longmapsto$} i^{*} : \quad \underline{\mathrm{Hom}}(A,\hand{$\mathcal{M}$}) \hand{$\longrightarrow$} \underline{\mathrm{Hom}}^{\hand{$!$}}(\hand{$\mathcal{M}$}, A^\wedge) \]
(which in any case is fully faithful) is equally an equivalence of catego-%
ries.

\margin{76}
